%%
%% Automatically generated file from DocOnce source
%% (https://github.com/doconce/doconce/)
%% doconce format latex week4.do.txt --latex_title_layout=beamer --latex_table_format=footnotesize --no_mako
%%
% #ifdef PTEX2TEX_EXPLANATION
%%
%% The file follows the ptex2tex extended LaTeX format, see
%% ptex2tex: https://code.google.com/p/ptex2tex/
%%
%% Run
%%      ptex2tex myfile
%% or
%%      doconce ptex2tex myfile
%%
%% to turn myfile.p.tex into an ordinary LaTeX file myfile.tex.
%% (The ptex2tex program: https://code.google.com/p/ptex2tex)
%% Many preprocess options can be added to ptex2tex or doconce ptex2tex
%%
%%      ptex2tex -DMINTED myfile
%%      doconce ptex2tex myfile envir=minted
%%
%% ptex2tex will typeset code environments according to a global or local
%% .ptex2tex.cfg configure file. doconce ptex2tex will typeset code
%% according to options on the command line (just type doconce ptex2tex to
%% see examples). If doconce ptex2tex has envir=minted, it enables the
%% minted style without needing -DMINTED.
% #endif

% #define PREAMBLE

% #ifdef PREAMBLE
%-------------------- begin preamble ----------------------

\documentclass[%
oneside,                 % oneside: electronic viewing, twoside: printing
final,                   % draft: marks overfull hboxes, figures with paths
10pt]{article}

\listfiles               %  print all files needed to compile this document

\usepackage{relsize,makeidx,color,setspace,amsmath,amsfonts,amssymb}
\usepackage[table]{xcolor}
\usepackage{bm,ltablex,microtype}

\usepackage[pdftex]{graphicx}

\usepackage{ptex2tex}
% #ifdef MINTED
\usepackage{minted}
\usemintedstyle{default}
% #endif

\usepackage[T1]{fontenc}
%\usepackage[latin1]{inputenc}
\usepackage{ucs}
\usepackage[utf8x]{inputenc}

\usepackage{lmodern}         % Latin Modern fonts derived from Computer Modern

% Hyperlinks in PDF:
\definecolor{linkcolor}{rgb}{0,0,0.4}
\usepackage{hyperref}
\hypersetup{
    breaklinks=true,
    colorlinks=true,
    linkcolor=linkcolor,
    urlcolor=linkcolor,
    citecolor=black,
    filecolor=black,
    %filecolor=blue,
    pdfmenubar=true,
    pdftoolbar=true,
    bookmarksdepth=3   % Uncomment (and tweak) for PDF bookmarks with more levels than the TOC
    }
%\hyperbaseurl{}   % hyperlinks are relative to this root

\setcounter{tocdepth}{2}  % levels in table of contents

% --- fancyhdr package for fancy headers ---
\usepackage{fancyhdr}
\fancyhf{} % sets both header and footer to nothing
\renewcommand{\headrulewidth}{0pt}
\fancyfoot[LE,RO]{\thepage}
% Ensure copyright on titlepage (article style) and chapter pages (book style)
\fancypagestyle{plain}{
  \fancyhf{}
  \fancyfoot[C]{{\footnotesize \copyright\ 1999-2026, Morten Hjorth-Jensen. Released under CC Attribution-NonCommercial 4.0 license}}
%  \renewcommand{\footrulewidth}{0mm}
  \renewcommand{\headrulewidth}{0mm}
}
% Ensure copyright on titlepages with \thispagestyle{empty}
\fancypagestyle{empty}{
  \fancyhf{}
  \fancyfoot[C]{{\footnotesize \copyright\ 1999-2026, Morten Hjorth-Jensen. Released under CC Attribution-NonCommercial 4.0 license}}
  \renewcommand{\footrulewidth}{0mm}
  \renewcommand{\headrulewidth}{0mm}
}

\pagestyle{fancy}


% prevent orhpans and widows
\clubpenalty = 10000
\widowpenalty = 10000

% --- end of standard preamble for documents ---


% insert custom LaTeX commands...

\raggedbottom
\makeindex

%-------------------- end preamble ----------------------

\begin{document}

% matching end for #ifdef PREAMBLE
% #endif

\newcommand{\exercisesection}[1]{\subsection*{#1}}


% ------------------- main content ----------------------



% ----------------- title -------------------------

\title{Quantum Computing, Quantum Machine Learning and Quantum Information Theories}

% ----------------- author(s) -------------------------

\author{Morten Hjorth-Jensen\inst{1}}
\institute{Department of Physics, University of Oslo, Oslo, Norway\inst{1}}
% ----------------- end author(s) -------------------------

\date{February 11, 2026
% <optional titlepage figure>
% <optional copyright>
}

% !split
\subsection{Plans for the week of February 9-13}

\begin{enumerate}
\item Reminder from last week on Entanglement, entropy and density matrices

\item One-qubit and two-qubit gates, background and realizations, see also PDF file URL:""

\item Simple Hamiltonian systems and start discussion of project 1

\item Introduction to qiskit (for exercise sessions).
% o \href{{https://youtu.be/UcfOVvFyw2E}}{Video of lecture}
% o \href{{https://github.com/CompPhysics/QuantumComputingMachineLearning/blob/gh-pages/doc/HandWrittenNotes/2025/NotesFebruary12.pdf}}{Whiteboard notes}
\end{enumerate}

\noindent
% !split
\subsection{Readings}

For the discussion of one-qubit, two-qubit and other gates, sections
2.6-2.11 and 3.1-3.4 of Hundt's book \textbf{Quantum Computing for Programmers},
contain most of the relevant information.

% !split
\subsection{Gates, the whys and hows}

In quantum computing it is common to rewrite various unitary
transformations acting on a given state, in terms of so-called gates
(one-qubit, two-qubit or more qubit gates). These unitary
transformations do actually represent specific interactions of the system with the
environment.

Each such operation is by convention written in terms of gates and a
chain of such gates represents a circuit. The latter represents then a
specific set of operations on an initial state in order to perform what we will label
as experiments.

The aim of the first set of notes this week is to link these gates
(and thereby circuits) to their respective unitary
transformation. These unitary transformations represent selected
physical processes. 

% !split
\subsection{Structure of the lecture}

\begin{enumerate}
\item First we review some of the basic ways of representing the solution to the Schrödinger equation, introducing the so-called Interaction, Heisenberg and Schrödinger prictures and unitary transformations. These material is meant mainly as background material

\item Secondly, we present examples of physical processes and how they can be represented as unitary operations on a given state.

\item These unitary transformations are then represented as gates. Setting gates together gives us a final circuit which can represent a specific physical system
\end{enumerate}

\noindent
% !split
\subsection{Part 1: Mathematical background}

The time-dependent Schrödinger equation (or equation of motion) reads
\[
\imath \hbar\frac{\partial }{\partial t}|\Psi_S(t)\rangle = \hat{H}\Psi_S(t)\rangle,
\]
where the subscript $S$ stands for Schrödinger here.
A formal solution is given by 
\[
|\Psi_S(t)\rangle = \exp{(-\imath\hat{H}(t-t_0)/\hbar)}|\Psi_S(t_0)\rangle.
\]

% !split
\subsection{Unitary transformation}

The Hamiltonian $\hat{H}$ is hermitian and the exponent represents a unitary 
operator with an operation carried out on the wave function at a time $t_0$.

The exponential term $\exp{(-\imath\hat{H}(t-t_0)/\hbar)}$ is our
unitary transformation $U(t,t_0)$ and we have
\[
|\Psi_S(t)\rangle = \exp{(-\imath\hat{H}(t-t_0)/\hbar)}|\Psi_S(t_0)\rangle=U(t,t_0)|\Psi_S(t_0)\rangle.
\]

% !split
\subsection{Interaction picture}

Our Hamiltonian is normally written out as the sum of an unperturbed part $\hat{H}_0$ and an interaction part $\hat{H}_I$, that is
\[
\hat{H}=\hat{H}_0+\hat{H}_I.
\]
In general we have $[\hat{H}_0,\hat{H}_I]\ne 0$ since $[\hat{T},\hat{V}]\ne 0$.
We wish now to define a unitary transformation in terms of $\hat{H}_0$ by defining
\[
|\Psi_I(t)\rangle = \exp{(\imath\hat{H}_0t/\hbar)}|\Psi_S(t)\rangle,
\]
which is again a unitary transformation carried out now at the time $t$ on the 
wave function in the Schrödinger picture. 

% !split
\subsection{Taking the derivative wrt time}

We can easily find the equation of motion by taking the time derivative
\[
\imath \hbar\frac{\partial }{\partial t}|\Psi_I(t)\rangle = -\hat{H}_0\exp{(\imath\hat{H}_0t/\hbar)}\Psi_S(t)\rangle+\exp{(\imath\hat{H}_0t/\hbar)}
\imath \hbar\frac{\partial }{\partial t}\Psi_S(t)\rangle.
\]
% !split
\subsection{Expression using the interaction picture}

Using the definition of the Schrödinger equation, we can rewrite the last equation as 
\[
\imath \hbar\frac{\partial }{\partial t}|\Psi_I(t)\rangle = \exp{(\imath\hat{H}_0t/\hbar)}\left[-\hat{H}_0+\hat{H}_0+\hat{H}_I\right]\exp{(-\imath\hat{H}_0t/\hbar)}\Psi_I(t)\rangle,
\]
which gives us
\[
\imath \hbar\frac{\partial }{\partial t}|\Psi_I(t)\rangle = \hat{H}_I(t)\Psi_I(t)\rangle,
\]
 with 
\[
\hat{H}_I(t)=
\exp{(\imath\hat{H}_0t/\hbar)}\hat{H}_I\exp{(-\imath\hat{H}_0t/\hbar)}.
\]

% !split
\subsection{Expectation value}

The order of the operators is important since $\hat{H}_0$ and $\hat{H}_I$ do generally not commute.
The expectation value of
an arbitrary operator in the interaction picture can now be written as
\[
\langle \Psi'_S(t)|\hat{O}_S|\Psi_S(t)\rangle = 
\langle \Psi'_I(t) |\exp{(\imath\hat{H}_0t/\hbar)}\hat{O}_I
\exp{(-\imath\hat{H}_0t/\hbar)}|\Psi_I(t)\rangle,
\]
and using the definition
\[
\hat{O}_I(t)=
\exp{(\imath\hat{H}_0t/\hbar)}\hat{O}_I\exp{(-\imath\hat{H}_0t/\hbar)},
\]
we obtain
\[
\langle \Psi'_S(t)|\hat{O}_S|\Psi_S(t)\rangle = 
\langle \Psi'_I(t) |\hat{O}_I(t)|\Psi_I(t)\rangle,
\]
stating that a unitary transformation does not change expectation values!

% !split
\subsection{Interaction picture and equation of motion}

If the take the time derivative of the operator in the interaction picture we arrive at the following equation of motion
\[
\imath \hbar\frac{\partial }{\partial t}\hat{O}_I(t) = \exp{(\imath\hat{H}_0t/\hbar)}\left[\hat{O}_S\hat{H}_0-\hat{H}_0\hat{O}_S\right]\exp{(-\imath\hat{H}_0t/\hbar)}=\left[\hat{O}_I(t),\hat{H}_0\right].
\]
Here we have used the time-independence of the Schrödinger equation
together with the observation that any function of an operator commutes with the operator itself. 

% !split
\subsection{Unitary operator}

In order to solve the equation of motion equation in the interaction picture, we define a unitary
time-development operator $\hat{U}(t,t')$. Later we will derive its
connection with the linked-diagram theorem, which yields a
linked expression for the actual operator. 
The action of the operator on the wave function is
\[
|\Psi_I(t) \rangle = \hat{U}(t,t_0)|\Psi_I(t_0)\rangle,
\]
with the obvious value $\hat{U}(t_0,t_0)=1$.

% !split
\subsection{Time-development operator}

The time-development operator $U$ has the
properties that
\[
     \hat{U}^{\dagger}(t,t')\hat{U}(t,t')=\hat{U}(t,t')\hat{U}^{\dagger}(t,t')=1,
\]
which implies that $U$ is unitary
\[
     \hat{U}^{\dagger}(t,t')=\hat{U}^{-1}(t,t').
\]
Further,
\[
    \hat{U}(t,t')\hat{U}(t't'')=\hat{U}(t,t'')
\]
and
\[
    \hat{U}(t,t')\hat{U}(t',t)=1,
\]
which leads to
\[
    \hat{U}(t,t')=\hat{U}^{\dagger}(t',t).
\]

% !split
\subsection{Properties of the operator $U$}

Using our definition of Schrödinger's equation in the interaction picture, we can then construct the operator $\hat{U}$. We have defined
\[
|\Psi_I(t)\rangle = \exp{(\imath\hat{H}_0t/\hbar)}|\Psi_S(t)\rangle,
\]
which can be rewritten as 
\[
|\Psi_I(t)\rangle = \exp{(\imath\hat{H}_0t/\hbar)}\exp{(-\imath\hat{H}(t-t_0)/\hbar)}|\Psi_S(t_0)\rangle,
\]
or
\[
|\Psi_I(t)\rangle = \exp{(\imath\hat{H}_0t/\hbar)}\exp{(-\imath\hat{H}(t-t_0)/\hbar)}\exp{(-\imath\hat{H}_0t_0/\hbar)}|\Psi_I(t_0)\rangle.
\]

% !split
\subsection{Interaction picture}

From the last expression we can define
\[
\hat{U}(t,t_0)=\exp{(\imath\hat{H}_0t/\hbar)}\exp{(-\imath\hat{H}(t-t_0)/\hbar)}\exp{(-\imath\hat{H}_0t_0/\hbar)}.
\]
It is then easy to convince oneself that the properties defined above are satisfied by the definition of $\hat{U}$. 

% !split
\subsection{Equation of motion}

We derive the equation of motion for $\hat{U}$ using the above definition.
This results in
\[
\imath \hbar\frac{\partial }{\partial t}\hat{U}(t,t_0) = \hat{H}_I(t)\hat{U}(t,t_0),
\]
which we integrate from $t_0$ to a time $t$ resulting in
\[
\hat{U}(t,t_0)-\hat{U}(t_0,t_0)=\hat{U}(t,t_0)-1=-\frac{\imath}{\hbar}\int_{t_0}^t dt' \hat{H}_I(t')\hat{U}(t',t_0),
\]
which can be rewritten as
\[
\hat{U}(t,t_0)=1-\frac{\imath}{\hbar}\int_{t_0}^t dt' \hat{H}_I(t')\hat{U}(t',t_0).
\]

% !split
\subsection{Time-ordering}

We can solve this equation iteratively keeping in mind the time-ordering of the operators
\[
\hat{U}(t,t_0)=1-\frac{\imath}{\hbar}\int_{t_0}^t dt' \hat{H}_I(t')+\left(\frac{-\imath}{\hbar}\right)^2\int_{t_0}^t dt'\int_{t_0}^{t'} dt'' \hat{H}_I(t')\hat{H}_I(t'')+\dots
\]
The third term can be written as 
\[
\int_{t_0}^t dt'\int_{t_0}^{t'} dt'' \hat{H}_I(t')\hat{H}_I(t'')=
\frac{1}{2}\int_{t_0}^t dt'\int_{t_0}^{t'} dt'' \hat{H}_I(t')\hat{H}_I(t'')
\]
\[
+\frac{1}{2}\int_{t_0}^t dt''\int_{t''}^{t} dt' \hat{H}_I(t')\hat{H}_I(t'').
\]

% !split
\subsection{Changing order of integrations}

We obtain this expression by changing the integration order in the second term
via a change of the integration variables $t'$ and $t''$  in 
\[
\frac{1}{2}\int_{t_0}^t dt'\int_{t_0}^{t'} dt'' \hat{H}_I(t')\hat{H}_I(t'').
\]
We can rewrite the terms which contain the double integral as
\[
\int_{t_0}^t dt'\int_{t_0}^{t'} dt'' \hat{H}_I(t')\hat{H}_I(t'')=
\]
\[
\frac{1}{2}\int_{t_0}^t dt'\int_{t_0}^{t'} dt''\left[\hat{H}_I(t')\hat{H}_I(t'')\Theta(t'-t'')
+\hat{H}_I(t')\hat{H}_I(t'')\Theta(t''-t')\right],
\]
with $\Theta(t''-t')$ being the standard Heavyside or step function. The step function allows us to give a specific time-ordering to the above expression.

% !split
\subsection{Time-ordering operator}

With the $\Theta$-function we can rewrite the last expression as 
\[
\int_{t_0}^t dt'\int_{t_0}^{t'} dt'' \hat{H}_I(t')\hat{H}_I(t'')=
\frac{1}{2}\int_{t_0}^t dt'\int_{t_0}^{t'} dt''\hat{T}\left[\hat{H}_I(t')\hat{H}_I(t'')\right],
\]
where $\Hat{T}$ is the so-called time-ordering operator. 

% !split
\subsection{Final expression for $U$}

With this definition, we can rewrite the expression for $\hat{U}$ as 
\[
\hat{U}(t,t_0)=\sum_{n=0}^{\infty}\left(\frac{-\imath}{\hbar}\right)^n\frac{1}{n!}
\int_{t_0}^t dt_1\dots \int_{t_0}^t dt_N \hat{T}\left[\hat{H}_I(t_1)\dots\hat{H}_I(t_n)\right]=
\]
\[
\hat{T}\exp{\left[\frac{-\imath}{\hbar}
\int_{t_0}^t dt' \hat{H}_I(t')\right]}.
\]

% !split
\subsection{Heisenberg  picture as alternative}

We wish now to define a unitary transformation in terms of $\hat{H}$ by defining
\[
|\Psi_H(t)\rangle = \exp{(\imath\hat{H}t/\hbar)}|\Psi_S(t)\rangle,
\]
which is again a unitary transformation carried out now at the time $t$ on the 
wave function in the Schrödinger picture. If we combine this equation with 
Schrödinger's equation we obtain the following equation of motion
\[
\imath \hbar\frac{\partial }{\partial t}|\Psi_H(t)\rangle = 0,
\]
meaning that $|\Psi_H(t)\rangle$ is time independent. An operator in this picture is defined as
\[
\hat{O}_H(t)=
\exp{(\imath\hat{H}t/\hbar)}\hat{O}_S\exp{(-\imath\hat{H}t/\hbar)}.
\]

% !split
\subsection{New equation of motion}

The time dependence is then in the operator itself, and this yields in turn the
following equation of motion
\[
\imath \hbar\frac{\partial }{\partial t}\hat{O}_H(t) = \exp{(\imath\hat{H}t/\hbar)}\left[\hat{O}_H\hat{H}-\hat{H}\hat{O}_H\right]\exp{(-\imath\hat{H}t/\hbar)}=\left[\hat{O}_H(t),\hat{H}\right].
\]
We note that an operator in the Heisenberg picture can be related to the corresponding
operator in the interaction picture as 
\[
\hat{O}_H(t)=
\exp{(\imath\hat{H}t/\hbar)}\hat{O}_S\exp{(-\imath\hat{H}t/\hbar)}=
\]
\[
\exp{(\imath\hat{H}_It/\hbar)}\exp{(-\imath\hat{H}_0t/\hbar)}\hat{O}_I\exp{(\imath\hat{H}_0t/\hbar)}\exp{(-\imath\hat{H}_It/\hbar)}.
\]

% !split
\subsection{Unitary transformations again}

With our definition of the time evolution operator we see that
\[
\hat{O}_H(t)=\hat{U}(0,t)\hat{O}_I\hat{U}(t,0),
\]
which in turn implies that $\hat{O}_S=\hat{O}_I(0)=\hat{O}_H(0)$, all operators are equal at $t=0$. The wave function in the Heisenberg formalism is 
related to the other pictures as 
\[
|\Psi_H\rangle=|\Psi_S(0)\rangle=|\Psi_I(0)\rangle,
\]
since the wave function in the Heisenberg picture is time independent. 
We can relate this wave function to that a given time $t$ via the time evolution operator as
\[
|\Psi_H\rangle=\hat{U}(0,t)|\Psi_I(t)\rangle.
\]

% !split
\subsection{Adiabatic switching}

We assume that the interaction term is switched on gradually. Our wave function at time $t=-\infty$ and $t=\infty$ is supposed to represent a non-interacting system
given by the solution to the unperturbed part of our Hamiltonian $\hat{H}_0$.
We assume the ground state is given by $|\Phi_0\rangle$, which could be a Slater determinant.
We define our Hamiltonian as
\[
\hat{H}=\hat{H}_0+\exp{(-\varepsilon t/\hbar)}\hat{H}_I,
\]
where $\varepsilon$ is a small number. The way we write the Hamiltonian 
and its interaction term is meant to simulate the switching of the interaction.

% !split
\subsection{Time evolution}

The time evolution of the wave function in the interaction picture is then
\[
|\Psi_I(t) \rangle = \hat{U}_{\varepsilon}(t,t_0)|\Psi_I(t_0)\rangle,
\]
with 
\[
\hat{U}_{\varepsilon}(t,t_0)=\sum_{n=0}^{\infty}\left(\frac{-\imath}{\hbar}\right)^n\frac{1}{n!}
\int_{t_0}^t dt_1\dots \int_{t_0}^t dt_N
\]
\[
\times \exp{(-\varepsilon(t_1+\dots+t_n)/\hbar)}\hat{T}\left[\hat{H}_I(t_1)\dots\hat{H}_I(t_n)\right]
\]

% !split
\subsection{Initial state preparation}

In the limit $t_0\rightarrow -\infty$, the solution ot Schrödinger's equation is
$|\Phi_0\rangle$, and the eigenenergies are given by 
\[
\hat{H}_0|\Phi_0\rangle=W_0|\Phi_0\rangle,
\]
meaning that 
\[
|\Psi_S(t_0)\rangle = \exp{(-\imath W_0t_0/\hbar)}|\Phi_0\rangle,
\]
with the corresponding interaction picture wave function given by
\[
|\Psi_I(t_0)\rangle = \exp{(\imath \hat{H}_0t_0/\hbar)}|\Psi_S(t_0)\rangle=|\Phi_0\rangle.
\]

% !split
\subsection{Final expression}

The solution becomes time independent in the limit $t_0\rightarrow -\infty$.
The same conclusion can be reached by looking at 
\[
\imath \hbar\frac{\partial }{\partial t}|\Psi_I(t)\rangle =
\exp{(\varepsilon |t|/\hbar)}\hat{H}_I|\Psi_I(t)\rangle 
\]
and taking the limit $t\rightarrow -\infty$.
We can rewrite the equation for the wave function at a time $t=0$ as
\[
|\Psi_I(0) \rangle = \hat{U}_{\varepsilon}(0,-\infty)|\Phi_0\rangle.
\]

% !split
\subsection{Part 2: Specific realizations and famous gates}

Nuclear magnetic resonance (NMR) quantum computing is one of the several
proposed approaches for constructing a quantum computer. It uses the
spin states of nuclei within molecules as qubits. The quantum states
are probed through the nuclear magnetic resonances, allowing the
system to be implemented as a variation of nuclear magnetic resonance
spectroscopy. NMR differs from other implementations of quantum
computers in that it uses an ensemble of systems, in this case
molecules, rather than a single pure state.

You can read more about this at \href{{https://cba.mit.edu/docs/papers/98.06.sciqc.pdf}}{\nolinkurl{https://cba.mit.edu/docs/papers/98.06.sciqc.pdf}}

% !split
\subsection{Spin Hamiltonian}

In order to understand in terms of a given Hamiltonian how the
different gates arise, we consider now the Hamiltonian of a nuclear
spin in a magnetic field. Since the spin provides provides a magnetic
dipole moment, a nucleus with a spin will interact with the magnetic
field. The Haniltonian of a nucleus with spin interacting with a
magnetic field $\bm{B}$ is

\[
H = -\bm{\mu}\bm{B},
\]
with $\bm{\mu}=\gamma\bm{S}$, $\gamma$ being the so-called gyromagnetic ratio and $\bm{S}$ the spin.

% !split
\subsection{Field along the $z$-axis}

It is common to let the spin interact with a constant magnetic field
along the $z$-axis. This gives an effecitve Hamiltonian

\[
H_z = -\frac{\hbar\omega_L}{2}\sigma_z,
\]

where $\omega_L$ is the so-called Larmor precession frequency. This
quantity includes also the constant magnetif field along the
$z$-axis. For all practical purposes it suffices for us to have an
expression of the Hamiltonian in terms of the Pauli-Z matrix.

% !split
\subsection{Bringing back a state on the Bloch sphere}

Suppose that our initial one qubit state (for example a spin-$1/2$
nucleus for NMR studies) points along some arbitrary axis. As
discussed during our second week, a point on the Bloch sphere can be
represented as at time $t=0$
\[
\vert \psi(t=0) \rangle = \vert \psi(0) \rangle=\cos{(\frac{\theta}{2})}\vert 0\rangle +\exp{\imath\phi}\sin{(\frac{\theta}{2})}\vert 1\rangle.
\]

% !split
\subsection{Time evolution}

Since the hamiltonian is time-independent, the state $\vert \psi(0)
\rangle$, our system will evolve according to the unitary transformation 
\[
\vert \psi(t) \rangle = U(t)\vert \psi(0) \rangle=\exp{\imath\omega_L t\sigma_z/2}\vert \psi(0) \rangle.
\]
Inserting the Bloch sphere ansatz we have then
\[
\vert \psi(t) \rangle=\exp{\imath\omega_L t\sigma_z/2}\cos{(\frac{\theta}{2})}\vert 0\rangle +\exp{\imath\omega_L t\sigma_z/2}\exp{\imath\phi}\sin{(\frac{\theta}{2})}\vert 1\rangle.
\]

The specific hamiltonian we have chosen here serves to exemplify how can represent physical operations in terms of specific gates, here a one-qubit gate (see whiteboard notes at \href{{https://github.com/CompPhysics/QuantumComputingMachineLearning/blob/gh-pages/doc/HandWrittenNotes/2025/NotesFebruary12.pdf}}{\nolinkurl{https://github.com/CompPhysics/QuantumComputingMachineLearning/blob/gh-pages/doc/HandWrittenNotes/2025/NotesFebruary12.pdf}}for more details).

% !split
\subsection{Final expression}
Assume we have a given operator $\bm{A}$ acting on a  vector space $\vert a\rangle$ with eigenvalues $a$ 
\[
\exp{\bm{A}}\vert a\rangle=\sum_{n=0}^{\infty} \frac{1}{n!}\bm{A}^n\vert a\rangle=\sum_{n=0}^{\infty} \frac{a^n}{n!}\vert a\rangle=\exp{a}\vert a\rangle.
\]

Using this result, we obtain
\[
\vert \psi(t) \rangle=\exp{\imath\omega_L t/2}\cos{(\frac{\theta}{2})}\vert 0\rangle +\exp{-\imath\omega_L t/2}\exp{\imath\phi}\sin{(\frac{\theta}{2})}\vert 1\rangle.
\]

% !split
\subsection{Part 3: Famous Quantum gates and quantum circuits}

Quantum gates are physical actions that are applied to the physical
system representing the qubits. Mathematically, they are
complex-valued, unitary matrices which act on the complex-values
normalized vectors that represent qubits. As the quantum analog of
classical logic gates (such as AND and OR), there is a corresponding
quantum gate for every classical gate; however, there are quantum
gates that have no classical counter-part. They act on a set of qubits
and, changing their state. That is, if $U$ is a quantum gate and
$\vert q \rangle $ is a qubit, then acting the gate $U$ on the qubit $\vert q \rangle $
transforms the qubit as follows:

\begin{align}
\vert q \rangle \overset{U}{\to}U\vert q \rangle 
.\end{align}

% !split
\subsection{Quantum circuits}

Quantum circuits are diagrammatic representations of quantum
algorithms. The horizontal dimension corresponds to time; moving left
to right corresponds to forward motion in time. They consist of a set
of qubits $\vert q_n\rangle$ which are stacked vertically on the left-hand
side of the diagram. Lines, called quantum wires, extend horizontally
to the right from each qubit, representing its state moving forward in
time. Additionally, they contain a set of quantum gates that are
applied to the quantum wires. Gates are applied chronologically, left
to right.

% !split
\subsection{Single-Qubit Gates}

A single-qubit gate is a physical action that is applied to one
qubit. It can be represented by a matrix $U$ from the group SU(2). Any
single-qubit gate can be parameterized by three angles: $\theta$,
$\phi$, and $\lambda$ as follows

\[
U(\theta,\phi,\lambda)=\begin{bmatrix}
\cos\frac{\theta}{2} & -e^{i\lambda}\sin\frac{\theta}{2}
\\
e^{i\phi}\sin\frac{\theta}{2} & e^{i(\phi+\lambda)}\cos\frac{\theta}{2}
\end{bmatrix}.
\]

% !split
\subsection{Widely used gates}

There are several widely used quantum gates. Perhaps the most famous are 
the Pauli gates correspond to the Pauli matrices

\[
I=\begin{bmatrix} 1 & 0 \\ 0 & 1 \end{bmatrix},
\]

\[
X =\begin{bmatrix} 0 & 1 \\ 1 & 0\end{bmatrix},
\]

\[
Y=\begin{bmatrix}0 & -i \\i & 0\end{bmatrix},
\]

\[
Z=\begin{bmatrix} 1 & 0 \\ 0 & -1\end{bmatrix}.
\]

% !split
\subsection{Algebra basis}

These gates form a basis for
the algebra $\mathfrak{su}(2)$. Exponentiating them will thus give us
a basis for SU(2), the group within which all single-qubit gates
live.

% !split
\subsection{Exponentiated Pauli gates}

These exponentiated Pauli gates are called rotation gates
$R_{\sigma}(\theta)$ because they rotate the quantum state around the
axis $\sigma=X,Y,Z$ of the Bloch sphere by an angle $\theta$. They are
defined as

\[
R_X(\theta)=e^{-i\frac{\theta}{2}X}=
\begin{bmatrix}
\cos\frac{\theta}{2} & -i\sin\frac{\theta}{2} \\
-i\sin\frac{\theta}{2} & \cos\frac{\theta}{2} 
\end{bmatrix},
\]
\[
R_Y(\theta)=e^{-i\frac{\theta}{2}Y}=
\begin{bmatrix}
\cos\frac{\theta}{2} & -\sin\frac{\theta}{2} \\
\sin\frac{\theta}{2} & \cos\frac{\theta}{2} 
\end{bmatrix},
\]
\[
R_Z(\theta)=e^{-i\frac{\theta}{2}Z}=\begin{bmatrix}
e^{-i\theta/2} & 0 \\
0 & e^{i\theta/2}\end{bmatrix}.
\]

% !split
\subsection{Basis for $\mathrm{SU}(2)$}

Because they form a basis for $\mathrm{SU}(2)$, any single-qubit gate
can be decomposed into three rotation gates. Indeed
\[
R_z(\phi)R_y(\theta)R_z(\lambda)=
\begin{bmatrix}
e^{-i\phi/2} & 0 \\
0 & e^{i\phi/2}
\end{bmatrix}
\begin{bmatrix}
\cos\frac{\theta}{2} & -\sin\frac{\theta}{2} \\
\sin\frac{\theta}{2} & \cos\frac{\theta}{2} 
\end{bmatrix}
\begin{bmatrix}
e^{-i\lambda/2} & 0 \\
0 & e^{i\lambda/2}
\end{bmatrix}
\]
which we can rewite as
\[
e^{-i(\phi+\lambda)/2}
\begin{bmatrix}
\cos\frac{\theta}{2} & -e^{i\lambda}\sin\frac{\theta}{2}\\
e^{i\phi}\sin\frac{\theta}{2} & e^{i(\phi+\lambda)}\cos\frac{\theta}{2}
\end{bmatrix},
\]

which is, up to a global phase, equal to the expression for an arbitrary single-qubit gate.

% !split
\subsection{Two-Qubit Gates}

A two-qubit gate is a physical action that is applied to two
qubits. It can be represented by a matrix $U$ from the group
SU(4). One important type of two-qubit gates are controlled gates,
which work as follows: Suppose $U$ is a single-qubit gate. A
controlled-$U$ gate ($CU$) acts on two qubits: a control qubit
$\vert x \rangle $ and a target qubit $\vert y \rangle $. The controlled-$U$ gate
applies the identity $I$ or the single-qubit gate $U$ to the target
qubit if the control gate is in the zero state $\vert 0\rangle$ or the one
state $\vert 1\rangle$, respectively.

% !split
\subsection{Control qubit}
The control qubit is not acted
upon. This can be represented as follows if
\[CU\vert xy\rangle=
\vert xy\rangle \hspace{0.1cm} \mathrm{if} \hspace{0.1cm}  \vert x \rangle =\vert 0\rangle.
\]

% !split
\subsection{In matrix form}

It is easier to see in a matrix form.
It can be written in matrix form by writing it as a superposition of
the two possible cases, each written as a simple tensor product

\[
CU = \vert 0\rangle\langle 0\vert\otimes I + \vert 1\rangle\langle 1 \vert \otimes U=\begin{bmatrix}
1 & 0 & 0 & 0 \\
0 & 1 & 0 & 0 \\
0 & 0 & u_{00} & u_{01} \\
0 & 0 & u_{10} & u_{11}
\end{bmatrix}.
\]

% !split
\subsection{CNOT gate}

One of the most fundamental controlled gates is the CNOT gate. It is
defined as the controlled-$X$ gate $CX$. It can be written in matrix form as follows:

\[
\mathrm{CNOT}=\mathrm{CX}=\begin{bmatrix}
1 & 0 & 0 & 0 \\
0 & 1 & 0 & 0 \\
0 & 0 & 0 & 1 \\
0 & 0 & 1 & 0
\end{bmatrix}.
\]

% !split
\subsection{$\mathrm{CX}$ gate}

It changes, when operating on a two-qubit state where the first qubit is the control qubit and the second qubit is the target qubit, the states (check this)
\[
\mathrm{CX}\vert 00\rangle=\vert 00\rangle,
\]
\[
\mathrm{CX}\vert 10\rangle= \vert 11\rangle,
\]
\[
\mathrm{CX}\vert 01\rangle= \vert 01\rangle,
\]
\[
\mathrm{CX}\vert 11\rangle= \vert 10\rangle,
\]
which you can easily see by simply multiplying the above matrix with any of the above states.

% !split
\subsection{Swap gate}

A widely used two-qubit gate that goes beyond the simple controlled function is the SWAP gate. It swaps the states of the two qubits it acts upon

\[
\mathrm{SWAP}\vert xy\rangle=\vert yx\rangle.
\]
and has the following matrix form

\[
\mathrm{SWAP}
=\begin{bmatrix}
1 & 0 & 0 & 0 \\
0 & 0 & 1 & 0 \\
0 & 1 & 0 & 0 \\
0 & 0 & 0 & 1
\end{bmatrix}.
\]

% !split
\subsection{Hamiltonians, one qubit system}

We start with a simple $2\times 2$ Hamiltonian matrix expressed in
terms of Pauli $X$ and $Z$ matrices, as discussed in the project text as well.

We define a  symmetric matrix  $H\in {\mathbb{R}}^{2\times 2}$
\[
H = \begin{bmatrix} H_{11} & H_{12} \\ H_{21} & H_{22}
\end{bmatrix},
\]

% !split
\subsection{Rewriting the Hamiltonian}

We  let $H = H_0 + H_I$, where
\[
H_0= \begin{bmatrix} E_1 & 0 \\ 0 & E_2\end{bmatrix},
\]
is a diagonal matrix. Similarly,
\[
H_I= \begin{bmatrix} V_{11} & V_{12} \\ V_{21} & V_{22}\end{bmatrix},
\]
where $V_{ij}$ represent various interaction matrix elements.
We can view $H_0$ as the non-interacting solution
\begin{equation}
       H_0\vert 0 \rangle =E_1\vert 0 \rangle,
\end{equation}
and
\begin{equation}
       H_0\vert 1\rangle =E_2\vert 1\rangle,
\end{equation}
where we have defined the orthogonal computational one-qubit basis states $\vert 0\rangle$ and $\vert 1\rangle$.

% !split
\subsection{Using Pauli matrices}

We rewrite $H$ (and $H_0$ and $H_I$)  via Pauli matrices
\[
H_0 = \mathcal{E} I + \Omega \sigma_z, \quad \mathcal{E} = \frac{E_1
  + E_2}{2}, \; \Omega = \frac{E_1-E_2}{2},
\]
and
\[
H_I = c \bm{I} +\omega_z\sigma_z + \omega_x\sigma_x,
\]
with $c = (V_{11}+V_{22})/2$, $\omega_z = (V_{11}-V_{22})/2$ and $\omega_x = V_{12}=V_{21}$.
We let our Hamiltonian depend linearly on a strength parameter $\lambda$

\[
H=H_0+\lambda H_\mathrm{I},
\]

with $\lambda \in [0,1]$, where the limits $\lambda=0$ and $\lambda=1$
represent the non-interacting (or unperturbed) and fully interacting
system, respectively.  The model is an eigenvalue problem with only
two available states.

% !split
\subsection{Selecting parameters}

Here we set the parameters $E_1=0$,
$E_2=4$, $V_{11}=-V_{22}=3$ and $V_{12}=V_{21}=0.2$.

The non-interacting solutions represent our computational basis.
Pertinent to our choice of parameters, is that at $\lambda\geq 2/3$,
the lowest eigenstate is dominated by $\vert 1\rangle$ while the upper
is $\vert 0 \rangle$. At $\lambda=1$ the $\vert 0 \rangle$ mixing of
the lowest eigenvalue is $1\%$ while for $\lambda\leq 2/3$ we have a
$\vert 0 \rangle$ component of more than $90\%$.  The character of the
eigenvectors has therefore been interchanged when passing $z=2/3$. The
value of the parameter $V_{12}$ represents the strength of the coupling
between the two states.

% !split
\subsection{Setting up the matrix}






















\bpycod
from  matplotlib import pyplot as plt
import numpy as np
dim = 2
Hamiltonian = np.zeros((dim,dim))
e0 = 0.0
e1 = 4.0
Xnondiag = 0.20
Xdiag = 3.0
Eigenvalue = np.zeros(dim)
# setting up the Hamiltonian
Hamiltonian[0,0] = Xdiag+e0
Hamiltonian[0,1] = Xnondiag
Hamiltonian[1,0] = Hamiltonian[0,1]
Hamiltonian[1,1] = e1-Xdiag
# diagonalize and obtain eigenvalues, not necessarily sorted
EigValues, EigVectors = np.linalg.eig(Hamiltonian)
permute = EigValues.argsort()
EigValues = EigValues[permute]
# print only the lowest eigenvalue
print(EigValues[0])

\epycod


Now rewrite it in terms of the identity matrix and the Pauli matrix X and Z

















\bpycod
# Now rewrite it in terms of the identity matrix and the Pauli matrix X and Z
X = np.array([[0,1],[1,0]])
Y = np.array([[0,-1j],[1j,0]])
Z = np.array([[1,0],[0,-1]])
# identity matrix
I = np.array([[1,0],[0,1]])

epsilon = (e0+e1)*0.5; omega = (e0-e1)*0.5
c = 0.0; omega_z=Xdiag; omega_x = Xnondiag
Hamiltonian = (epsilon+c)*I+(omega_z+omega)*Z+omega_x*X
EigValues, EigVectors = np.linalg.eig(Hamiltonian)
permute = EigValues.argsort()
EigValues = EigValues[permute]
# print only the lowest eigenvalue
print(EigValues[0])

\epycod


% !split
\subsection{Where are we going?}

For a one-qubit system we can reach every point on the Bloch sphere
(as discussed earlier) with a rotation about the $x$-axis and the
$y$-axis.

We can express this mathematically through the following operations (see whiteboard for the drawing), giving us a new state $\vert \psi\rangle$
\[
\vert\psi\rangle = R_y(\phi)R_x(\theta)\vert 0 \rangle.
\]

% !split
\subsection{Possible ansatzes}

We can produce multiple ansatzes for the new state in terms of the
angles $\theta$ and $\phi$.  With these ansatzes we can in turn
calculate the expectation value of the above Hamiltonian, now
rewritten in terms of various Pauli matrices (and thereby gates), that is compute

\[
\langle \psi \vert (c+\mathcal{E})\bm{I} + (\Omega+\omega_z)\bm{\sigma}_z + \omega_x\bm{\sigma}_x\vert \psi \rangle.
\]

We can now set up a series of ansatzes for $\vert \psi \rangle$ as
function of the angles $\theta$ and $\phi$ and find thereafter the
variational minimum using for example a gradient descent method.

% !split
\subsection{More on rotation operators}

To do so, we need to remind ourselves about the mathematical expressions for
the rotational matrices/operators.

\[
R_x(\theta)=\cos{\frac{\theta}{2}}\bm{I}-\imath \sin{\frac{\theta}{2}}\bm{\sigma}_x,
\]

and

\[
R_y(\phi)=\cos{\frac{\phi}{2}}\bm{I}-\imath \sin{\frac{\phi}{2}}\bm{\sigma}_y.
\]

% !split
\subsection{Code example}

















\bpycod
# define the rotation matrices
# Define angles theta and phi
theta = 0.5*np.pi; phi = 0.2*np.pi
Rx = np.cos(theta*0.5)*I-1j*np.sin(theta*0.5)*X
Ry = np.cos(phi*0.5)*I-1j*np.sin(phi*0.5)*Y
#define basis states
basis0 = np.array([1,0])
basis1 = np.array([0,1])

NewBasis = Ry @ Rx @ basis0
print(NewBasis)
# Compute the expectation value
#Note hermitian conjugation
Energy = NewBasis.conj().T @ Hamiltonian @ NewBasis
print(Energy)

\epycod

Not an impressive results. We set up now a loop over many angles $\theta$ and $\phi$ and compute the energies


















\bpycod
# define a number of angles
n = 20
angle = np.arange(0,180,10)
n = np.size(angle)
ExpectationValues = np.zeros((n,n))
for i in range (n):
    theta = np.pi*angle[i]/180.0
    Rx = np.cos(theta*0.5)*I-1j*np.sin(theta*0.5)*X
    for j in range (n):
        phi = np.pi*angle[j]/180.0
        Ry = np.cos(phi*0.5)*I-1j*np.sin(phi*0.5)*Y
        NewBasis = Ry @ Rx @ basis0
        Energy = NewBasis.conj().T @ Hamiltonian @ NewBasis
        Edifference=abs(np.real(EigValues[0]-Energy))
        ExpectationValues[i,j]=Edifference

print(np.min(ExpectationValues))

\epycod


Clearly, this is not the very best way of proceeding. Rather, here we
would compute the gradient and thereby find the minimum as function of
the angles $\theta$ and $\phi$.  This will lead us to the so-called Variational Quantum Eigensolver to be discussed next week.

% !split
\subsection{Plans for week 5, February 16-20}

\begin{enumerate}
\item We will extend the above one-qubit Hamiltonian to a two-qubit problem and analyze how we can set up its simulation

\item Thereafter we will focus on the variational quantum eigensolver for simulating the Hamiltonian systems discussed in exercises a-d) of project 1.

\item For the exercises next week, we will focus on discussions of parts a-d) of project 1. See also the exercises here.
\end{enumerate}

\noindent
% !split
\subsection{Exercises and tasks for next week: Hamiltonians and project 1}

As an initial test, we consider a simple $2\times 2$ real
Hamiltonian consisting of a diagonal part $H_0$ and off-diagonal part
$H_I$, playing the roles of a non-interactive one-body and interactive
two-body part respectively. Defined through their matrix elements, we
express them in the Pauli basis $\vert 0\rangle$ and $\vert 1 \rangle$

\begin{align*}
    \begin{split} 
        H &= H_0 + H_I \\
        H_0 = \begin{bmatrix}
            E_1 & 0 \\
            0 & E_2
        \end{bmatrix}&, \hspace{20px}
        H_I = \lambda \begin{bmatrix}
            V_{11} & V_{12} \\
            V_{21} & V_{22}
        \end{bmatrix}
    \end{split}
\end{align*}
Where $\lambda \in [0,1]$ is a coupling constant parameterizing the strength of the interaction. 

% !split
\subsection{Rewriting in terms of Pauli matrices}

Define
\[
    E_+ = \frac{E_1 + E_2}{2},\hspace{20px} E_- = \frac{E_1 - E_2}{2}
\]
show  that by combining the identity and $Z$ Pauli matrix, this can be expressed as

\[
    H_0 = E_+ I + E_- Z
\]

% !split
\subsection{The interaction part}

For $H_1$ we use the same trick to fill the diagonal, defining $V_+ = (V_{11} + V_{22})/2, V_- = (V_{11} - V_{22})/2$. From the hermiticity requirements of $H$, we note that $V_{12} = V_{21} \equiv V_o$. Use this to simplify the problem to a simple $X$ term. 

\[
    H_I = V_+ I + V_- Z + V_o X
\]

% !split
\subsection{Measurement basis}

For the above system show that the Pauli $X$ matrix can be rewritten in terms of the Hadamard matrices and the Pauli $Z$ matrix, that is
\[
X=HZH.
\]


% ------------------- end of main content ---------------

% #ifdef PREAMBLE
\end{document}
% #endif

